% =====================================================================
% Results — skeleton, 2026-08-12.
%
% Structure is frozen; values marked \NUM{} are placeholders to be filled from
% the artifact named in the comment above each block once the runs land.
%
% The organising principle: not every experiment is here to say "latency wins".
% Sections A and B establish the champion. Section C explains *why* it is a good
% operating point. Section D delimits where the advantage stops. A paper whose
% every result points the same direction reads as advocacy; this one should not.
% =====================================================================

\section{Results and Discussion}

\subsection{Design-space screening: the top-ranked configuration}
\label{sec:res-rank}

% results/analysis/v1/{mean_rank.csv,friedman.json}  ->  paper/tables/tab_ranking_v1.tex
Across all $27$ configurations evaluated over $20$ paired (dataset $\times$ seed)
blocks, \texttt{LeakyParallel/latency} attains the lowest mean macro-F1 rank
($4.03$). A Friedman test rejects the null of equal performance
($\chi^2=263.9$, $p=3.2\times10^{-41}$), and the Nemenyi critical difference is
$9.28$.

That critical difference is the reason we describe this configuration as
\emph{top-ranked} rather than as statistically best: $14$ of the $27$
configurations lie within one critical difference of the leader, so the ranking
identifies a leading group and \texttt{LeakyParallel/latency} sits at the head
of it. In particular \texttt{Leaky/latency} is statistically indistinguishable
from it ($p=0.33$), and we treat the two as co-leaders on terminal quality.

% results/analysis/v1/focused_tests.json -> paper/tables/tab_focused_v1.tex
The pre-specified comparison against the closest competitor,
\texttt{LeakyParallel/rate}, favours latency in $16$ of $20$ blocks with a
median difference of $+0.70$ pp macro-F1 (Holm-adjusted $p=0.0084$). The
hierarchical bootstrap over datasets gives $+3.79$ pp $[+0.11, +10.19]$. The
interval is wide because the effect is strongly dataset-dependent: it is
$+13.3$ pp on CIC-IDS2017 and near zero on the other three. At the encoding
level, aggregating over all nine neurons, the block-level test favours latency
over rate (Holm $p=0.027$) while the dataset-level bootstrap interval
$[-0.18, +11.30]$ includes zero. We report both. Taken together they describe an
effect that is consistent in direction but variable in magnitude, which is a
weaker and more accurate claim than uniform superiority.

We note one point of interpretation, because it is easy to reach the opposite
conclusion from the wrong metric. On KDDCup99, rate marginally exceeds latency
in \emph{accuracy} ($0.9260$ vs $0.9257$); on macro-F1, the pre-registered
primary metric for these heavily imbalanced datasets, latency leads by
$0.379$ pp. Fixing the primary metric before the analysis is what makes that
resolution credible rather than convenient.

\subsection{Strict confirmation along both design axes}
\label{sec:res-strict}

% results/analysis/v2_encoding/, results/analysis/v2_latency/
Under the leakage-resistant protocol we confirm the selection along each axis
separately. With the neuron fixed, latency remains the strongest encoding on
\NUM{n} of the evaluated protocols (Table~\NUM{tab}). With the encoding fixed,
\texttt{LeakyParallel} remains among the highest-ranked neuron families
(Table~\NUM{tab}). We report per-dataset means with across-seed standard
deviations throughout; no dispersion in this paper is pooled across datasets.

% results/v1_submitted/ctu13_strict_vanilla.csv
The clearest effect of the strict protocol is on CTU-13, and it is large enough
that the original number should not stand. Under the submitted protocol ---
whole-scenario host aggregates with a random flow split ---
\texttt{LeakyParallel/latency} attains a perfect $100.00\pm0.00$ macro-F1.
Making the aggregation causal alone leaves $99.69\pm0.23$. Adding a
scenario-disjoint partition brings it to $\mathbf{92.01\pm9.72}$, with folds at
$99.9$, $99.9$, $91.2$ and $77.0$. We therefore report $92.01$ with its fold
range as the CTU-13 result, and retain the perfect score only as a labelled
protocol comparison. The fold spread is itself informative: it reflects genuine
differences in scenario difficulty that a uniform perfect score concealed
entirely, and it is the reason we never quote the mean alone.

\subsection{Mechanism: latency concentrates class evidence earlier}
\label{sec:res-readout}

% results/analysis/early_readout/summary.csv, figures readout_*.pdf
Terminal macro-F1 alone would leave the encoding recommendation resting on a
modest and dataset-dependent margin. The anytime readout gives a mechanistic
account instead. On \NUM{protocol}, latency coding reaches $99\%$ of its
terminal macro-F1 at a median of $T_{99}=3$ timesteps, against $T_{99}=20$ for
rate. Truncating the readout at $T_{99}$ discards $74.2\%$ of latency's output
spikes but only $24.2\%$ of rate's.

The timing advantage is not uniform across initialisations and we do not present
it as such: latency's $T_{99}$ has an interquartile range of $2$--$9$ and a full
range of $2$--$23$ across five seeds, against $20$--$21$ and $17$--$23$ for rate.
The distributions are largely disjoint but one latency seed converges no faster
than rate. Figure~\NUM{fig} plots every seed individually rather than a mean
band for this reason.

% results/analysis/efficiency/summary.csv
The synaptic-operation accounting explains the mechanism directly. Latency emits
$14.6$ input spikes per sample --- one per feature, by construction --- while
rate's Bernoulli draw emits $238.6$ over the same window; with an input fanout
of $128$ the input layer dominates the operation count. The result is
$3{,}363$ SOPs per sample for latency against $32{,}502$ for rate, a factor of
$9.7$, at \emph{higher} terminal quality. Latency is therefore Pareto-optimal
against rate on quality, decision latency and operation count simultaneously.

This matters for how much weight the delta result has to carry. Delta reaches
$99\%$ of its terminal score at $T_{99}=1$, but that terminal score is
$6.5$ pp below latency's: it stabilises early on a worse answer, which is not an
early-inference advantage. Consistent with our pre-registration, the encoding
conclusion rests on the latency-versus-rate comparison, and delta's weakness is
reported as a property of applying a temporal-contrast code to static per-flow
vectors rather than as evidence for latency.

\subsection{Where the advantage does not extend}
\label{sec:res-limits}

Three results bound the claim, and we report them because a configuration
recommendation is only useful alongside its operating envelope.

% results/analysis/security/fixed_far.csv
\textbf{Operating thresholds do not transfer under distribution shift.}
Selecting a threshold on validation to hold false alarms at $1\%$ and freezing
it yields roughly $7.7\%$ on the official NSL-KDD test set. KDDTest+ is
deliberately distribution-shifted relative to KDDTrain+, so a threshold
calibrated on data drawn from the training distribution does not hold. This is
an unfavourable operational number and a substantive one: it says that
reporting a test-selected ROC operating point --- as is common --- would
overstate deployable performance.

% results/analysis/known_unseen/known_unseen.csv
\textbf{Unseen attack types remain hard, and latency does not help there.}
$29.2\%$ of KDDTest+ attack rows belong to $17$ types absent from training.
Detection falls from $0.8023$ on known types to $0.3481$ on unseen types for
\texttt{LeakyParallel/latency}. Notably latency has the \emph{best} known-type
detection of the three encodings and the \emph{worst} unseen-type detection. We
therefore use this experiment to delimit the recommendation rather than to
support it: \texttt{LeakyParallel/latency} is the strongest configuration for
the supervised detection task we evaluate, and that advantage does not extend to
attack types absent from training.

% results/analysis/security/duplicate_free.csv
\textbf{Benchmark duplication is large but its effect is small.} $57.53\%$ of
the official KDDCup99 test set occurs verbatim in training. Restricting
evaluation to the non-duplicated remainder costs $0.79$ pp macro-F1. The
overlap is severe enough that absolute KDDCup99 numbers should not be read as
evidence of generalisation, but it does not account for the ranking.

\subsection{Summary}

Across a controlled $27$-configuration \texttt{snnTorch} design-space
evaluation, \texttt{LeakyParallel} with latency encoding attains the lowest mean
macro-F1 rank, with a leading group of configurations statistically
indistinguishable from it. Focused leakage-resistant experiments confirm the
selection along both design axes, and the anytime-readout analysis shows that
latency coding concentrates discriminative evidence substantially earlier than
rate coding, yielding a Pareto-optimal operating point in quality, decision
latency and synaptic operations. Strict evaluation simultaneously lowers the
optimistic estimates the original protocol produced --- most sharply on CTU-13
--- and exposes two limits of the recommendation: calibrated operating
thresholds do not transfer under distribution shift, and unseen attack types
remain difficult.
